% Report in Springer LNCS format
% Compile with:
%   pdflatex kg_proposal_validation.tex && pdflatex kg_proposal_validation.tex

\documentclass{llncs}

\usepackage[utf8]{inputenc}
\usepackage[T1]{fontenc}
\usepackage{microtype}
\usepackage{amsmath,amssymb}
\usepackage{booktabs}

\begin{document}

\title{Link Prediction on FB15k-237:\\
Validation of Self-Adversarial Negative Sampling}

\author{Thalia Delhaise \and Lenny Briclet \and
        Andr\'e Fonseca \and Rafael Gufler}

\institute{RC2526, Knowledge Graphs}

\maketitle

% -----------------------------------------------------------------------
\begin{abstract}
We study link prediction on the FB15k-237 benchmark using two knowledge
graph embedding models, TransE and RotatE, trained under matched
conditions. Our current baseline shows that RotatE outperforms TransE with
default Bernoulli negative sampling. We now propose to validate
self-adversarial negative sampling as the main extension of the project.
The goal is to test whether focusing on harder negatives improves global
performance and whether the gains are concentrated on difficult graph
structures such as rare relations and low-degree entities.
\end{abstract}

% -----------------------------------------------------------------------
\section{Project Goal}

We study link prediction on the FB15k-237 knowledge graph benchmark using
two standard knowledge graph embedding models: TransE and RotatE.
Our goal is to compare both models under matched training conditions and
then study whether better negative sampling improves performance on the
harder parts of the graph.

Our current baseline results show that RotatE outperforms TransE under
the same setup:

\begin{table}[h]
\centering
\caption{Current test-set results with default Bernoulli negative sampling.}
\label{tab:baseline}
\begin{tabular}{@{}lcccc@{}}
\toprule
Model & MRR & Hits@1 & Hits@3 & Hits@10 \\
\midrule
TransE & 0.152 & 0.073 & 0.174 & 0.304 \\
RotatE & 0.261 & 0.179 & 0.286 & 0.424 \\
\bottomrule
\end{tabular}
\end{table}

These results were obtained with matched settings: embedding dimension
$d=128$, Adam optimizer with learning rate $10^{-3}$, batch size 1024,
up to 50 epochs, early stopping on validation MRR, CUDA, and seed 42.

% -----------------------------------------------------------------------
\section{Proposed Extension: Self-Adversarial Negative Sampling}

During training, negative triples are created by corrupting true triples.
For example, from a true triple $(h,r,t)$, we can replace the tail entity
and create a fake triple $(h,r,t')$.

The issue with standard random or Bernoulli negative sampling is that many
negative triples are too easy: they are obviously false, so after a few
epochs the model learns little from them. We therefore propose to use
\textbf{self-adversarial negative sampling}, introduced with RotatE, which
gives more importance to negative triples that the current model still
finds plausible.

Our hypothesis is that this will not improve all cases uniformly. Instead,
we expect the largest gains on harder parts of the graph, especially
low-frequency relations and low-degree entities, where random negatives
are often less informative.

For each positive triple $(h,r,t)$, we sample a pool of $n$ candidate
negative triples. In the tail-corruption case:

\[
\{(h,r,t'_1), (h,r,t'_2), \ldots, (h,r,t'_n)\}.
\]

Let $s_\theta(h,r,t)$ be a score where higher means more plausible.
Each negative receives the weight:

\begin{equation}
  p_j
  =
  \frac{
    \exp\!\bigl(\alpha s_\theta(h,r,t'_j)\bigr)
  }{
    \sum_{i=1}^{n}
    \exp\!\bigl(\alpha s_\theta(h,r,t'_i)\bigr)
  },
  \label{eq:sadv}
\end{equation}

where $\alpha > 0$ controls how strongly the method focuses on hard
negatives.

The loss is:

\begin{equation}
  \mathcal{L}
  =
  -\log\sigma\!\bigl(s_\theta(h,r,t)\bigr)
  -
  \sum_{j=1}^{n}
  p_j
  \log\sigma\!\bigl(-s_\theta(h,r,t'_j)\bigr).
  \label{eq:loss}
\end{equation}

The first term increases the score of the true triple, while the second
term decreases the scores of the weighted negative triples.

In the main experiment, we plan to use $n=64$ candidate negatives. If this
is too expensive computationally, we will reduce it to $n=32$ and report
this choice explicitly. We will tune
$\alpha \in \{0.5, 1.0, 2.0\}$ using the validation set only.

% -----------------------------------------------------------------------
\section{Experimental Design}

The main experiment is a clean ablation:

\begin{center}
\textbf{RotatE with default Bernoulli negatives}
\quad vs. \quad
\textbf{RotatE with self-adversarial negatives}.
\end{center}

Everything else will remain fixed: same dataset, same model architecture,
same embedding size, same optimizer, same batch size, same training
budget, same seed, and same filtered evaluation protocol.

If time allows, we will also run the same comparison on TransE. However,
the minimum required experiment is the RotatE ablation, since the method
was introduced for RotatE and RotatE is our stronger baseline.

We will select $\alpha$ using validation MRR and then report final
filtered MRR and Hits@10 on the test set.

% -----------------------------------------------------------------------
\section{Planned Analysis}

In addition to global MRR and Hits@10, we will perform slice analysis to
understand where the method helps.

We will report results separately for:

\begin{itemize}
  \item \textbf{relation-frequency slices}: rare vs. frequent relations;
  \item \textbf{entity-degree slices}: low-degree vs. high-degree entities;
  \item \textbf{head vs. tail prediction}: because the baseline already
        shows a large asymmetry between head and tail queries.
\end{itemize}

All relation frequencies and entity degrees will be computed from the
training set only, to avoid using test-set information during analysis.

This will allow us to check whether self-adversarial negative sampling
gives a general improvement or mainly helps on specific difficult parts of
the graph.

% -----------------------------------------------------------------------
\section{Questions for Feedback}

We would like feedback on the following points before implementing the
method:

\begin{itemize}
  \item Is self-adversarial negative sampling an appropriate extension for
        the scope of this project?
  \item Is the ablation design sound, given that we keep all settings fixed
        and only change the negative sampler?
  \item Is fixing the candidate pool size to $n=64$ and tuning only
        $\alpha \in \{0.5,1.0,2.0\}$ reasonable?
  \item Is the planned slice analysis sufficient to support our main claim
        about where hard negatives help?
  \item Would you recommend testing another negative sampling strategy
        instead, or is this one suitable?
\end{itemize}

\end{document}